\chapter{Type Expansion, Ideation, and Agents}

The repository's prompt-driven agent is not an accidental extra on top of the typing
system. It is one of the clearest demonstrations that the same semantic stance can govern
planning and generation as well as checking.

\section*{Type expansion}
The type-expansion modules in the hybrid package take a short natural-language description
and elaborate it into a much richer plan: modules, covers, dependencies, obligations,
interfaces, and eventually code-generation structure. This process is easiest to interpret
as growth in a presheaf. The intent layer starts with sparse, natural-language local
sections. Structural, semantic, proof, and code layers gradually acquire more content as
the plan becomes more explicit.

The important conceptual shift is that generation is not described as free-form text
completion followed by a verifier. It is described as typed elaboration. The prompt gives
an intent section. Expansion grows a cover. Local generation produces candidate sections.
Verification checks whether they glue. Repair loops try to resolve obstructions. Assembly
produces a global artifact.

\section*{Ideation as typing}
The ideation machinery is even more unusual. The repository includes a large mathematical
ontology and domain-selection logic intended to enrich prompts by analogy to concepts from
many mathematical areas. The right way to phrase this in the paper is not that the system
has a magical creativity engine. It is that it tries to treat ideation itself as a typed
search over richer local descriptions.

If a prompt about trading is enriched by analogy to algebraic geometry, order theory,
signal processing, or topology, the resulting design search space is larger but also more
structured. Ideation becomes a pre-typing phase that improves the quality of the later
cover and the plausibility of the resulting local sections.

\section*{Scalar-valued typing}
The ideation pipeline can be sharpened still further by introducing what the repository now
calls \emph{scalar-valued typing}. The idea is simple: before committing to a full
elaboration, assign a scalar to the candidate elaboration that measures how strongly the
current domain evidence supports it. In implementation terms, \repo{agent/cli.py} already
tracks an LLM belief score, a mathematically-derived score based on validated analogies and
certifiable structure, and per-area strengths. Scalar-valued typing packages these into a
single typed quantity attached to the elaboration.

If $E$ is a type elaboration and $d$ is a target domain, then the scalar-valued typing
score can be written schematically as
\[
  \Score(E,d) =
  \alpha \, \Score_{\mathrm{LLM}}(E,d)
  + (1-\alpha)\, \Score_{\mathrm{Math}}(E,d)
  + \varepsilon \, \max_i \Score_i(E,d),
\]
with \(\alpha = 0.55\) and \(\varepsilon = 0.05\) in the current implementation. The
first term measures how strongly the language model believes the idea explains the problem
in the chosen domain. The second term measures structural richness: validated analogies,
deep connections, theorem schemas, certification targets, and penalties for unresolved
\(H^1\)-mass. The final term lets especially strong local mathematical areas gently boost a
globally promising elaboration.

This is a useful addition because it makes ideation behave less like unranked brainstorming
and more like typed search. Scalar-valued typing does not replace ordinary type checking.
Instead it orders candidate elaborations before the expensive downstream steps of full
module expansion, proof planning, and code generation. The score is therefore neither a
proof nor a mere UX flourish; it is a typed prioritization signal.

\begin{figure}[h]
\centering
\begin{tikzpicture}[
  node distance=1.15cm,
  box/.style={draw=chapterblue, rounded corners=2pt, fill=chapterblue!6, minimum width=3.1cm, minimum height=0.9cm, align=center},
  sum/.style={draw=accentgold, rounded corners=2pt, fill=accentgold!12, minimum width=3.5cm, minimum height=0.95cm, align=center},
  >=Latex
]
\node[box] (intent) {prompt / intent section};
\node[box, right=of intent] (llm) {LLM belief \\ \(\Score_{\mathrm{LLM}}\)};
\node[box, right=of llm] (math) {math-model evidence \\ \(\Score_{\mathrm{Math}}\)};
\node[box, below=of llm] (areas) {area strengths \\ \(\max_i \Score_i\)};
\node[sum, below=of math] (overall) {scalar-valued typing \\ \(\Score(E,d)\)};
\draw[->, thick, chapterblue] (intent) -- (llm);
\draw[->, thick, chapterblue] (intent.east) .. controls +(1.5,-0.4) and +(-1.4,0.6) .. (math.west);
\draw[->, thick, chapterblue] (llm) -- (overall);
\draw[->, thick, chapterblue] (math) -- (overall);
\draw[->, thick, chapterblue] (areas) -- (overall);
\end{tikzpicture}
\caption{Scalar-valued typing adds a numeric layer over ideation: type elaborations are scored by domain fit before full expansion.}
\end{figure}

\section*{The orchestration agent}
The agent orchestrator in \repo{agent/orchestrator.py} embodies a multi-stage state
machine: parsing, expanding, planning, generating, verifying, and assembling. The code is
therefore easier to understand after one adopts the hybrid type picture. Each stage is not
simply passing strings to the next. It is refining what kind of section exists, which
layer it inhabits, and which trust-bearing obligations must be checked before the next
stage is allowed to treat the result as global enough.

\begin{algorithm}[h]
\caption{Cohomological generation loop}
\begin{algorithmic}[1]
\State Resolve user prompt into intent sections.
\State Expand intent into a typed cover of local module obligations.
\For{each local module site}
  \State Synthesize candidate code and local contracts.
  \State Check structural, semantic, and proof obligations.
  \If{an obstruction remains}
    \State refine the local section and retry
  \EndIf
\EndFor
\State Attempt to glue the verified local sections into a project.
\State Emit residual obstructions if gluing fails.
\end{algorithmic}
\end{algorithm}

The importance of this loop is that it places generation inside the same semantics as
checking. A paper that really integrates the repository should emphasize this strongly. In
the current repository story, scalar-valued typing sits just before the expensive center of
this loop, biasing the search toward elaborations whose local mathematical evidence most
strongly supports the target domain.
